\selectlanguage{ngerman}%
\frenchspacing%
\begin{ledgroupsized}[r]{120mm}%
\footnotesize%
\pstart%
\noindent\textbf{Überlieferung:}%
\pend%
\end{ledgroupsized}%
\begin{ledgroupsized}[r]{114mm}%
\footnotesize%
\pstart%
\parindent -6mm%
\makebox[6mm][l]{\textit{L}}%
Auszüge mit Bemerkungen aus
\textsc{C.\:F.\:M.\;Dechales},%
\indexN{\textls{Dechales} (Chalesius), Claude François Milliet 1621\textendash1678}
\textit{Cursus seu mundus mathematicus} (3~Bde, Lyon 1674)%
:\cite{00124}
LH~XXXV\;9,\;17 Bl.\;5--6.
Ein Bogen 2\textsuperscript{o}.
Zweizweidrittel Seiten auf Bl.\;5\,r\textsuperscript{o} bis Bl.\;6\,r\textsuperscript{o},
zumeist einspaltig beschrieben.
Bl.\;6\,r\textsuperscript{o} überliefert auch N.\;\ref{41182-2} (Auszüge zur Statik aus demselben Werk) in diesem Band.
Bl.\;6\,v\textsuperscript{o} ist unbeschrieben.
Der Text ist durch Abstände und Striche unterteilt;
die Anordnung der %
Auszüge entspricht nicht durchweg der Textabfolge in der %
Vorlage.
\textemdash\
Weitere Auszüge aus Dechales' \textit{Cursus} überliefern im Leibniz-Nachlass die Manuskripte:
LH~XXXV\;9,\;17 Bl.\;1--4; %
LH~XXXVI Bl.\;54--57, 60--61; %
LH~XXXVII\;5 Bl.\;2--3 %
(sie werden voraussichtlich in künftigen Bänden von \textit{LSB} VIII ediert).
Ein Exemplar von Dechales' \textit{Cursus} mit Gebrauchsspuren von Leibnizens Hand liegt in der GWLB Hannover unter der Signatur Nm-A\,10018:1--3 vor.
\pend%
\end{ledgroupsized}%
\vspace{5mm}%
\begin{ledgroup}%
\footnotesize%
\pstart%
\noindent%
\textbf{Datierungsgründe:}
Die Datierung der Auszüge mit Bemerkungen N.\;\ref{41182-1} beruht auf Eigenschaften des Textträgers sowie auf inhaltlichen Betrachtungen.
Siehe die ausführliche Begründung auf S.\;\pageref{41182_dat}.%
\pend%
\end{ledgroup}%
\selectlanguage{latin}%
\frenchspacing%
\vspace{8mm}%
\normalsize%
\pstart%
\noindent%
\centering%
\bracks{%
\href{\Scan{LH_35_09_17_005r}}%
{5\,r\textsuperscript{o}}%
}
\hfill%
\textls{Statica}%
\indexS{statica}
\hfill%
\phantom{\bracks{5\,r\textsuperscript{o}}}%
\pend%
\vspace{0.5em}%
\pstart%
\noindent%
Incertum an in locis centro terrae%
\indexS{centrum terrae}
propioribus%
\indexS{locus propior centro terrae}
minor gravitatio gravium.%
\indexS{gravitatio gravium}%
\indexS{grave in loco propiori centro terrae} %
\pend%
\pstart%
\edtext{%
\textit{Corpus non possumus}%
\indexS{corpus separatum a terra}
separare a terra%
\indexS{terra}
\textit{nisi instrumentis}%
\indexS{istrumentum annexum terrae}
ipsi \textit{terrae annexis.
Aviculae}%
\indexS{avicula}
cum \textit{in aera tolluntur
primo saltant pedibus%
\indexS{pes aviculae}
terramque percutiunt,
deinde alis%
\indexS{ala aviculae}
aerem contra terram}%
\indexS{aer impulsus contra terram}
\edtext{%
\textit{impellunt}
(+\;nil%
}{%
\lemma{\textit{impellunt}}%
\Bfootnote{%
\bars{, deinde alis aerem contra terram impellunt \textit{streicht Hrsg.}} (+\;nil%
~\textit{L}%
}}
refert an contra terram%
\;\bracks{+})
et aer ipse incumbit terrae,%
\indexS{aer incumbens terrae}
ut aqua.%
\indexS{aqua incumbens terrae}%
}{%
\lemma{\textit{Corpus} \bracks{...} aqua}%
\Cfootnote{%
Vgl. \textsc{C.\,F.\,M.\;Dechales},
\textit{Cursus seu mundus mathematicus}, tract.\;VII, lib.\;I, digr.\;III
(Lyon 1674, Bd.\;I, S.\;434b:\cite{00124}
Zitat mit Auslassungen).%
}}
\edtext{%
Venti quoque%
\indexS{ventus separans gravia a terra}
cum gravia a terra separantur
sunt motus aeris,%
\indexS{ventus ut motus aeris}%
\indexS{motus aeris}
et ipsorum motus etiam
\edtext{%
oritur quodam nisu repellendi terram.%
\indexS{nisus repellendi terram}%
\indexS{terra}%
}{%
\lemma{oritur}\Bfootnote{%
\textit{(1)}~a corpore quod terrae %
\textit{(2)}~quodam nisu repellendi terram.%
~\textit{L}%
}}%
}{%
\lemma{Venti \bracks{...} terram}%
\Cfootnote{%
Vgl. a.a.O. (S.\;434b--435a).\cite{00124}%
}}
\pend%
\pstart%
\edtext{%
\textit{Si aer esset levis,
condensatus} foret \textit{levior},%
\indexS{aer condensatus}
quod contra est.%
}{%
\lemma{\textit{Si aer} \bracks{...} contra est}%
\Cfootnote{%
Vgl. a.a.O., digr.\;V
(S.\;437a--b).\cite{00124}%
}}
\pend%
\pstart%
\edtext{%
\textls{Thermometrum} simplex%
\indexS{thermometrum simplex}
\edtext{%
\textit{A} vel \textit{B}%
}{%
\lemma{\textit{A} vel \textit{B}}%
\Cfootnote{%
Vgl. die Diagramme \bracks{\textit{Fig.\:1}} u.\;\bracks{\textit{Fig.\:2}}.%
}}
nempe calefit
et inde
\edtext{%
in aqua%
}{%
\lemma{in}%
\Bfootnote{%
\textit{(1)}~aerem %
\textit{(2)}~aqua%
~\textit{L}%
}}
osculo%
\indexS{osculum thermometri}
\edtext{%
aperto \textit{C} collocatur,%
}{%
\lemma{aperto\;\textit{C}}%
\Bfootnote{%
\textit{(1)}~vel %
\textit{(2)}~collocatur,%
~\textit{L}%
}}
refrigefactum attrahet aerem,%
\indexS{thermometrum refrigefactum attrahens aerem}
itaque calefactum rursus expellet,%
\indexS{thermometrum calefactum expellens aerem}
vel saltem depellet.%
}{%
\lemma{\textls{Thermometrum} \bracks{...} depellet}%
\Cfootnote{%
Vgl. a.a.O. (S.\;441b--442a).\cite{00124}%
}}
\pend%
\newpage%
\begin{edfigure}[0.00em][2.50em]%
\hfill
\includegraphics[scale=0.85]{%
gesamttex/edit_VIII-4/images/LH_35_09_17_005-006_d01_005r.pdf}
\hfill
\includegraphics[scale=0.85]{%
gesamttex/edit_VIII-4/images/LH_35_09_17_005-006_d02_005r.pdf}
\hfill
\includegraphics[scale=0.85]{%
gesamttex/edit_VIII-4/images/LH_35_09_17_005-006_d03_005r.pdf}
\hfill
\includegraphics[scale=0.85]{%
gesamttex/edit_VIII-4/images/LH_35_09_17_005-006_d04_005r.pdf}
\hfill
\edcaption{%
\bracks{\textit{Fig.\:1}}%
\hspace*{20mm}%
\bracks{\textit{Fig.\:2}}%
\hspace*{20mm}
\bracks{\textit{Fig.\:3}}%
\hspace*{20mm}%
\bracks{\textit{Fig.\:4}}%
\edtext{}{%
\lemma{\bracks{\textit{Fig.\:1}} u. \bracks{\textit{Fig.\:2}}}%
\Cfootnote[nonum]{%
Vgl. \textsc{Dechales}, \textit{Cursus}, tract.\;VII, lib.\;I, %
digr.\;VI (S.\;441b: einzige Diagramme).\cite{00124}%
}}%
\edtext{}{%
\lemma{\bracks{\textit{Fig.\:3}} u. \bracks{\textit{Fig.\:4}}}%
\Cfootnote[nonum]{%
Vgl. a.a.O. (S.\;442a: erstes u. zweites Diagramm).\cite{00124}%
}}%
}
\end{edfigure}
\pstart%
\edtext{%
Alia
\edtext{%
species thermometri \textit{HGF},%
\indexS{species thermometri}%
\indexS{thermometrum alterius speciei}%
}{%
\lemma{species thermometri \textit{HGF}}%
\Cfootnote{%
Vgl. das Diagramm \bracks{\textit{Fig.\:3}}.%
}}
ubi frigore%
\indexS{frigus contrahens aerem}
aer \textit{GF} contrahitur,%
\indexS{aer contractus frigore}
et aqua descendit,%
\indexS{aqua ascendens in thermometro}
cum in priori descendat calore.%
\indexS{calor}%
}{%
\lemma{Alia \bracks{...} calore}%
\Cfootnote{%
Vgl. a.a.O. (S.\;442a).\cite{00124}%
}}
\pend%
\pstart%
\edtext{%
\textit{Tertia%
\indexS{thermometrum tertiae speciei}
\edtext{%
species \textit{KMN},%
}{%
\lemma{\textit{species \textit{KMN}}}%
\Cfootnote{%
Vgl. das Diagramm \bracks{\textit{Fig.\:4}}.%
}}
aqua \textit{KM} descendet calore%
\indexS{aqua descendens in thermometro}
ex \textit{M}}
ad \textit{R},
\textit{ascendet ex \textit{K}}
ad \textit{O}.
Addi potest index%
\indexS{index in trochlea}
qui in
\edtext{%
trochlea \textit{S}%
\indexS{trochlea in thermometro}%
}{%
\lemma{trochlea \textit{S}}%
\Cfootnote{%
Vgl. das Diagramm \bracks{\textit{Fig.\:5}}.%
}}
notet numeris
quantum ascenderit \textit{KO},
globulus autem ligneus%
\indexS{globulus ligneus natans in aqua thermometri}
in \textit{K} natat in aqua.%
\indexS{aqua thermometri}%
}{%
\lemma{\textit{Tertia} \bracks{...} aqua}%
\Cfootnote{%
Vgl. \textsc{Dechales}, \textit{Cursus}, tract.\;VII, lib.\;I, %
digr.\;VI (S.\;442a).\cite{00124}%
}}
\pend%
\begin{edfigure}[2.50em][0.00em]%
\hfill%
\includegraphics[scale=0.85]{%
gesamttex/edit_VIII-4/images/LH_35_09_17_005-006_d05_005r.pdf}
\hfill%
\edcaption{%
\bracks{\textit{Fig.\:5}}
\edtext{}{%
\lemma{\bracks{\textit{Fig.\:5}}}%
\Cfootnote[nonum]{%
Vgl. \textsc{Dechales}, \textit{Cursus}, tract.\;VII, lib.\;I, %
digr.\;VI (S.\;442a: viertes Diagramm).\cite{00124}
Ein gestrichener Entwurf wird nicht wiedergegeben.}}
}
\end{edfigure}%
\newpage%
\pstart%
\edtext{%
\textit{Globulus vitreus%
\indexS{globulus vitreus mergens frigore}
inferiori parte apertus
semiplenus aere,%
\edtext{}{%
\lemma{\textit{Am Rand:}}%
\Afootnote{NB}}
ingruente frigore%
\indexS{frigus ingruens}
magis mergitur,}
tempore calido minus.%
\indexS{tempus calidum}%
}{%
\lemma{\textit{Globulus} \bracks{...} minus}%
\Cfootnote{%
Vgl. a.a.O. (S.\;442a).\cite{00124}%
}}
\edlabel{LH_35_09_17_005-006_005r_ewk-1}
\edtext{%
Huc pertinent \textls{\textit{cucurbitulae}}%
\indexS{cucurbitula medicorum}
\edtext{%
\textit{A},%
}{%
\lemma{\textit{A}}%
\Bfootnote{%
\textit{erg.\:L}%
}}
\textit{calefit aer inclusus quantum potest,%
\indexS{aer inclusus cucurbitulae}%
\indexS{aer cucurbitulae calefactus}
deinde} quantum potest
\textit{carni applicatur firmiter,%
\indexS{cucurbitula applicata carni}
ut nullus circa ejus labra%
\indexS{labrum cucurbitulae}
aeri externo pateat aditus,%
\indexS{aditus patens aeri externo}%
\indexS{aer externus}
tum frigefit}%
\indexS{cucurbitula frigefacta}
\textit{vel} af\textit{fusa frigid}\angles{\textit{a}}%
\indexS{frigida (aqua) affusa}
\textit{vel pannis madidis applicatis},%
\indexS{pannus madidus}
et \textit{caro elevatur intra cucurbitulam,%
\indexS{caro elevata intra cucurbitulam}
et saepe etiam ossa
si quae sint luxata.%
\indexS{os luxatum}
Si incidatur caro%
\indexS{caro incisa}
multus sanguis exilit.%
\indexS{sanguis exiliens}}%
\edlabel{LH_35_09_17_005-006_005r_ewk-2}%
}{%
\lemma{Huc \bracks{...} \textit{exilit}}%
\Cfootnote{%
Vgl. a.a.O. (S.\;442a:\cite{00124}
Zitat mit Auslassungen).
Die Passage ist zweimal exzerpiert;
vgl. unten, \refpassage{LH_35_09_17_005-006_005v_dfg-1}{LH_35_09_17_005-006_005v_dfg-2}.%
}}
\pend%
\begin{edfigure}[0.5em][1.0em]%
\includegraphics[scale=0.90]{%
gesamttex/edit_VIII-4/images/LH_35_09_17_005-006_d06_005r.pdf}
\edcaption[0.25em]{%
\hspace{4mm}%
\bracks{\textit{Fig.\:6}}%
\edtext{}{%
\lemma{\bracks{\textit{Fig.\:6}}}%
\Cfootnote[nonum]{%
Vgl. \textsc{Dechales}, \textit{Cursus}, tract.\;VII, lib.\;I, %
digr.\;VI (S.\;442a: fünftes Diagramm).\cite{00124}%
}}%
}%
\end{edfigure}
\pstart%
\edtext{%
\textit{R.\:P.\;Franc.\;de la Chaise}%
\indexN{\textls{La Chaise} (La Chaize, Hedraeus), François de 1624\textendash1709}
sumsit \textit{tubum 56 pedum%
\indexS{pes}
secundum parietis altitudinem erectum%
\indexS{tubus erectus 56 pedum}
aqua plenum,
ubi orificium
intra aliam aquam} demersum
\textit{fuit apertum}%
\bracks{,}
\textit{mansit tantum aqua}
\edtext{%
\bracks{\textit{ad}}%
}{%
\lemma{\textit{ad}}%
\Bfootnote{%
\textit{erg. Hrsg. nach Vorlage}%
}}
\textit{altitudinem pedum 32 circiter,%
\indexS{aqua 32 pedum altitudinis}
qui aequales in gravitate%
\indexS{gravitas 32 pedum aquae}%
\indexS{gravitas 2,25 pedum mercurii}
pedibus duobus et~%
$\displaystyle\frac{\textit{1}}{\textit{4}}$} Mercurii.%
\indexS{mercurius 2,25 pedum altitudinis}%
}{%
\lemma{\textit{R.\:P.\;Franc.} \bracks{...} Mercurii}%
\Cfootnote{%
Vgl. a.a.O., digr.\:VII, Nr.\;14
(S.\;443b--444a).\cite{00124}
François de La Chaise (1624--1709)%
\indexN{\textls{La Chaise} (La Chaize, Hedraeus), François de 1624\textendash1709}
war Jesuit und Beichtvater Ludwigs XIV.%
\indexN{\textls{Frankreich: Ludwig XIV.}, König 1643\textendash1715}
Dechales%
\indexN{\textls{Dechales} (Chalesius), Claude François Milliet 1621\textendash1678}
beschreibt ihn als \textit{in hujusmodi experimentis physicis exercitatissimus} (ebd.).%
}}
\edlabel{LH_35_09_17_005-006_005r_svj-1}%
\edtext{%
Si tubus
Mercurium%
\indexS{mercurius suspensus in tubo}
\edtext{%
suspensum%
}{%
\lemma{suspensum}%
\Bfootnote{%
\textit{erg.\;L}%
}}
continens%
\indexS{tubus continens mercurium sustentatus e bilance}
ex bilance
\edtext{%
sustentatur,
aequilibrium habebitur pondere%
\indexS{aequilibrium tubi continentis mercurium}%
}{%
\lemma{sustentatur,}%
\Bfootnote{%
\textit{(1)}~pondere %
\textit{(2)}~aequilibrium habebitur pondere%
~\textit{L}%
}}
quod et tubo et mercurio simul sit aequale.%
\edlabel{LH_35_09_17_005-006_005r_svj-2}%
}{%
\lemma{Si tubus \bracks{...} aequale}%
\Cfootnote{%
Vgl. a.a.O., N.\;16 (S.\;444a).\cite{00124}
Der Auszug wird später fortgesetzt;
vgl. unten \refpassage{LH_35_09_17_005-006_005v_ukp-1}{LH_35_09_17_005-006_005v_ukp-2}
u. das Diagramm \bracks{\textit{Fig.\:11}}.%
}}
\pend%
\pstart%
\edtext{%
\textit{Pascalius%
\indexN{\textls{Pascal} (Pascalius), Blaise 1623\textendash1662}
in Arvernis\indexO{Auvergne}
prope Clermontium}\indexO{Clermont-Ferrand}
hoc curavit fieri experimentum%
\indexS{experimentum Pascalii}%
\bracks{:}\;%
\textit{ad radices altissimi montis%
\indexS{radix montis}%
\indexS{mons altissimus}
le puits de dome}\indexO{Puy de Dôme}
\edtext{%
mercurius \textit{supra superficiem
stagnantis in vase} alterius mercurii%
\indexS{mercurius stagnans}
attollebatur%
}{%
\lemma{mercurius}%
\Bfootnote{%
\textit{(1)}~attollebatur %
\textit{(2)}~\textit{supra superficiem} %
\bracks{...} mercurii attollebatur%
~\textit{L}%
}}
\textit{digitis 26 et 4 lineis}%
\bracks{;}
conscenso monte \textit{ad altitudinem}
\textit{perpendicularem hexapedarum 150} %
\textit{altitudo}
\textit{intra tubum fuit digitorum 25,%
\indexS{altitudo mercurii intra tubum}%
\indexS{mercurius intra tubum}
\ledpb%
et superato montis vertice%
\indexS{vertex montis}
ad 500 hexapedas,%
\indexS{hexapeda}
digitorum 23 et}
\edtext{%
\bracks{\textit{duarum}}%
}{%
\lemma{trium}%
\Bfootnote{%
\textit{L\;ändert Hrsg. nach Vorlage}%
}}
\textit{linearum}%
\bracks{:}
\textit{quare altitudo diminuta
tribus digitis%
\indexS{digitus}
et 2 lineis.%
\indexS{linea (mensura)}}%
}{%
\lemma{\textit{Pascalius} \bracks{...} \textit{lineis}}%
\Cfootnote{%
Vgl. \textsc{Dechales}, \textit{Cursus}, tract.\:VII, lib.\;I, %
digr.\:VII, Nr.\;18 (S.\;444a:\cite{00124}
Zitat mit Auslassungen).
Siehe zum Experiment am Puy de Dôme\indexO{Puy de Dôme} etwa:
B.\;\textsc{Pascal},
\textit{Récit de la grande expérience}, 1648\cite{02577} %
(\textit{Traitez}, Paris 1663, S.\;165--194;\cite{00081}
\textit{PO}\;II, S.\;653--690).\cite{01524}
}}
\pend%
\pstart%
\edtext{%
\edlabel{LH_35_09_17_005-006_005r_scs-1}%
\textit{Qui Mercurium} habent \textit{in officinis},%
\indexS{officina}%
\indexS{mercurius servatus in officinis}
dicunt
\textit{quantacunque} eum tractent \textit{diligentia},%
\indexS{diligentia in mercurio servando}
tamen quovis
\edtext{%
anno
\textit{20\textsuperscript{ma} aut 25\textsuperscript{ta} parte decrescere.}%
\edlabel{LH_35_09_17_005-006_005r_scs-2}%
}{%
\lemma{anno}%
\Bfootnote{%
\textit{(1)}~perdere %
\textit{(2)}~\textit{20\textsuperscript{ma} aut 25\textsuperscript{ta} parte decrescere.}%
~\textit{L}%
}}
Ergo non mirum
si inde vapores sublati ferantur%
\indexS{vapor mercurii in vacuitate tubi}
in vacuitatem tubi,%
\indexS{vacuitas tubi}
et ibi dilatentur.%
}{%
\lemma{\textit{Qui} \bracks{...} dilatentur}%
\Cfootnote{%
Vgl. \textsc{Dechales}, \textit{Cursus}, tract.\:VII, lib.\;I, %
digr.\:VII (S.\;445a--b).\cite{00124}
Die Passage ist zum Teil zweimal exzerpiert;
vgl. unten,
\refpassage{LH_35_09_17_005-006_005v_jkr-1}{LH_35_09_17_005-006_005v_jkr-2}.%
}}
\pend%
\vspace{1.0em}%
\pstart%
\noindent%
\edtext{%
\textit{%
Lampades accensae et carbones igniti%
\indexS{lampas in recipiente evacuato}%
\indexS{carbo in recipiente evacuato}
citius extinguuntur in recipiente%
\indexS{recipiens evacuatus}
quam in aere libero,%
\indexS{aer liber}
et adhuc citius}
cum \textit{exugitur recipientis aer.%
\indexS{aer exuctus e recipiente}
Funis militaris accensus%
\indexS{funis militaris}
et in recipiens demissus
primo fumo implevit recipiens,}%
\indexS{recipiens impletus fumo}
exucto vero aere cum fumo,%
\indexS{fumus exuctus cum aere}%
\indexS{aer exuctus cum fumo}
\textit{ignis non apparuit,}%
\indexS{ignis in recipiente evacuato}%
\indexS{aer immissus recipienti}
sed \textit{immisso} iterum \textit{aere
coepit} denuo \textit{ardere.}%
}{%
\lemma{\textit{Lampades} \bracks{...} \textit{ardere}%
}%
\Cfootnote{%
Vgl. \textsc{Dechales}, \textit{Cursus}, tract.\:VII, lib.\;I, %
digr.\;IX, experim.\;17
(S.\;455b: Zitat mit Auslassungen);\cite{00124}
nach \textsc{R.\;Boyle}, \textit{New experiments}, Nr.\;10--12 %
(Oxford 1660, S.\;74--83;\cite{00015}
\textit{BW}~I, S.\;184--188).\cite{00017}%
}}
(+\;%
Experiendum
quamdiu hoc modo
in
\edtext{%
vacuo vis%
}{%
\lemma{vacuo}%
\Bfootnote{%
\textit{(1)}~actio %
\textit{(2)}~vis~%
\textit{L}%
}}
ignis sopita%
\indexS{vis ignis sopita}%
\indexS{vis ignis conservata in vacuo}%
\indexS{vacuum}
conservari possit.%
~\bracks{+)}
\pend%
\pstart%
\edtext{%
Elastrum%
\indexS{elastrum sclopeti}
\edtext{%
sclopeti%
\indexS{sclopetum}%
}{%
\lemma{sclopeti}%
\Bfootnote{%
\textit{erg.\;L}%
}}
intra recipientem exhaustum%
\indexS{recipiens exhaustus}
egit ut extra,
imo \textit{scintillas excussit}%
\indexS{scintilla}
quibus \textit{accensus est pulvis} pyrius,%
\indexS{pulvis pyrius}
\textit{flammamque%
\indexS{flamma pulveris pyrii}
habuit more consueto.}%
\indexS{mos consuetus}%
}{%
\lemma{Elastrum \bracks{...} \textit{consueto}}%
\Cfootnote{%
Vgl. \textsc{Dechales}, \textit{Cursus}, tract.\:VII, lib.\;I, digr.\;IX, %
experim.\;18 (S.\;455b: Zitat mit Auslassungen);\cite{00124}
nach \textsc{Boyle}, \textit{New experiments}, Nr.\;14
(S.\;88--102;\cite{00015} \textit{BW}~I, S.\;189--191).\cite{00017}%
}}
\pend%
\pstart%
\edtext{%
\textit{Ex oleo%
\indexS{oleum}
erumpunt plures bullae,%
\indexS{bulla erumpens}
ex aqua distillata%
\indexS{aqua distillata}
pauciores
quam ex communi.%
\indexS{aqua communis}%
} %
Aqua in recipiente posita%
\indexS{aqua in recipiente exhausto}%
\indexS{recipiens exhaustus}
in \textit{phiala colli oblongi},%
\indexS{phiala colli oblongi}
exucto aere%
\indexS{aer exuctus e recipiente}
nonnihil intumuit,
cum prius \textit{bullae%
\indexS{bulla erumpens}
a fundo ad summitatem usque erupissent.}%
}{%
\lemma{\textit{Ex oleo} \bracks{...} erupissent}%
\Cfootnote{%
Vgl. \textsc{Dechales}, \textit{Cursus}, tract.\:VII, lib.\;I, digr.\;IX, %
experim.\;22 (S.\;456b:\cite{00124}
Zitat mit Auslassungen);\cite{00124}
nach \textsc{Boyle}, \textit{New experiments}, Nr.\;23--24 %
(S.\;182--189;\cite{00015} \textit{BW}~I, S.\;221--223).\cite{00017}%
}}
\pend%
\vspace{1.0em}
\begin{edfigure}[0.5em][0.0em]%
\includegraphics[scale=0.85]{%
gesamttex/edit_VIII-4/images/LH_35_09_17_005-006_d07_005v.pdf}%
\edcaption[0.25em]{%
\bracks{\textit{Fig.\:7}}%
\edtext{}{%
\lemma{\lbrack\textit{Fig.\:7}\rbrack}%
\Cfootnote[nonum]{%
Vgl. \textsc{Dechales}, \textit{Cursus}, tract.\:VII, lib.\;I, digr.\;IX, %
experim.\;5 (S.\;450a: drittes Diagramm).\cite{00124}%
}}%
}%
\end{edfigure}%
\newpage%
\pstart%
\noindent%
\edtext{%
\edtext{%
\textit{LM} vas aqua plenum%
\indexS{vas plenum aqua}%
}{%
\lemma{\textit{LM}}%
\Bfootnote{%
\textit{(1)}~aqua plena %
\textit{(2)}~\bars{\textls{vas} \textit{erg.}} aqua plenum%
~\textit{L}%
}}
usque ad osculum,%
\indexS{osculum vasis}
in eo natat globulus \textit{N} %
semiplenus aqua%
\indexS{globulus natans semiplenus aqua}%
\bracks{;}
si vasculi orificium%
\indexS{orificium vasculi}
et ita aquam digitis%
\indexS{digitus}
premas%
\bracks{,}
aer in \textit{N} contentus comprimetur%
\indexS{aer compressus}%
\bracks{;}
ergo plus aquae ingredietur,
et mergi poterit \textit{N}.
\edtext{%
Pro meliori compressione%
\indexS{compressio aeris}
quidam orificio%
\indexS{orificium vasculi}
\textit{circumligant vesicam suillam.}%
\indexS{vesica suilla}%
}{%
\lemma{Pro meliori}%
\Bfootnote{%
\bracks{...} \textit{vesicam suillam.} %
\textit{erg.\;L}
}}
Quidam \textit{artificium} ita \textit{occultant,%
\indexS{artificium occultatum}
suilla vesica collo addita%
\indexS{vesica suilla}
totam phialam%
\indexS{phiala inversa}
invertunt,}
et foramen inserunt mensae,%
\indexS{foramen in mensa}%
\indexS{mensa forata}
\textit{ita ut vesica suilla%
\indexS{vesica suilla}
plena aqua intra mensam depend}\angles{\textit{eat}},
\textit{tum lignum ita accommodant%
\indexS{lignum accommodatum}
ut adduci et remitti
\edtext{%
possit,
et per%
}{%
\lemma{\textit{possit, et}}%
\Bfootnote{%
\bars{et \textit{streicht Hrsg.}} \textit{per}%
~\textit{L}%
}}
adductionem vesicam comprimat.%
\indexS{vesica suilla}%
}
Quidam
\edtext{%
\textit{extra cubiculum}%
\indexS{cubiculus}%
}{%
\lemma{\textit{extra}}%
\Bfootnote{%
\textit{(1)}~vesicam %
\textit{(2)}~\textit{cubiculum}%
~\textit{L}%
}}
\textit{positus,
ampullis imperabat,%
\indexS{ampulla ascendens vel descendens}
ut ascenderent vel descenderent,}
colore%
\indexS{color}
autem distinguebantur
et \textit{filum adducebat ad certam mensuram.%
\indexS{filum adductum ad mensuram}%
}%
}{%
\lemma{\textit{LM} \textls{vas} \bracks{...} \textit{mensuram}}%
\Cfootnote{%
Vgl. a.a.O. (S.\;450a: Zitat mit Auslassungen).\cite{00124}
Die Beschreibung bezieht sich auf das Diagramm \bracks{\textit{Fig.\:7}.}%
}}
\bracks{%
\href{\Scan{LH_35_09_17_005v}}%
{5\,v\textsuperscript{o}}%
}
\pend%
\pstart%
Elegantissime res
praestaretur hoc
\edtext{%
modo%
\bracks{:}
vasi vitreo \textit{L}%
\indexS{vas vitreum}
qualicunque,%
}{%
\lemma{modo}%
\Bfootnote{%
\textit{(1)}~vas\angles{e} %
\textit{(2)}~vas vitreum \bars{\textit{L} \textit{erg.}} lagen %
\textit{(3)}~vasi vitreo \textit{L} qualicunque,%
~\textit{L}%
}}
in quo natat
\edlabel{35_09_17_005-006_005v_tbe-1a}%
ampulla~\textit{N},%
\indexS{ampulla natans in vase}%
\edtext{}{%
\xxref{35_09_17_005-006_005v_tbe-1a}{35_09_17_005-006_005v_tbe-1b}%
\lemma{ampulla~\textit{N} \bracks{...} in~\textit{N}}%
\Cfootnote{%
Die Bezeichnung \textit{N} wird im Text sowie im entsprechenden Diagramm \bracks{\textit{Fig.\:8}} zweideutig verwendet.%
}}
superponatur aliud vas vitreum inversum majus \textit{MNP}.%
\indexS{vas vitreum}
\edtext{%
In mensa%
\indexS{mensa forata}
cui%
}{%
\lemma{In}%
\Bfootnote{%
\textit{(1)}~fundo cui %
\textit{(2)}~mensa cui%
~\textit{L}%
}}
insistunt
sit foramen \textit{Q}%
\indexS{foramen in mensa}
per quod aer%
\indexS{aer inflatus ope tubi}
ope tubi inflatus%
\indexS{tubus ad aerem inflandum}
insensibiliter aquam comprimet,%
\indexS{aer comprimens aquam}%
\indexS{aqua compressa aere}
et aerem in~\textit{N}.%
\indexS{aer compressus}%
\edlabel{35_09_17_005-006_005v_tbe-1b}
Foramen \textit{Q} sit insensibile%
\bracks{,}
vas inversum \textit{MNP}%
\indexS{vas vitreum}
sit sub jugo \textit{RS},%
\indexS{jugum}
ne attollatur flatu,%
\indexS{flatus aeris}
vel alia ratione firmetur.%
\indexS{ratio firmitatis}
Potest quis efficere
ut quaelibet ampulla%
\indexS{ampulla natans in vase}
maneat in quolibet \angles{loco}.
Patietur enim ascendere 
\edtext{%
aliquousque
vel descendere,%
\indexS{ampulla ascendens vel descendens}
et ibi%
}{%
\lemma{aliquousque}%
\Bfootnote{%
\textit{(1)}~, et ibi %
\textit{(2)}~vel descendere, et ibi%
~\textit{L}%
}}
laxabit.%
\edtext{}{%
\lemma{\textit{Am Ende der Zeile:}}%
\Afootnote{%
\Denarius%
}}
\pend%
\begin{edfigure}%
\hspace{-55mm}%
\includegraphics[scale=0.95]{%
gesamttex/edit_VIII-4/images/LH_35_09_17_005-006_d08_005r.pdf}
\edcaption{%
\hspace{-55mm}%
\lbrack\textit{Fig.\:8}\rbrack}
\vspace{-6.50em}          
\hspace{55mm}%
\includegraphics[scale=0.95]{%
gesamttex/edit_VIII-4/images/LH_35_09_17_005-006_d09_005v.pdf}
\edcaption[1.00em]{%
\hspace{55mm}%
\lbrack\textit{Fig.\:9}\rbrack%
\edtext{}{%
\lemma{\lbrack\textit{Fig.\:9}\rbrack}%
\Cfootnote[nonum]{%
Vgl. \textsc{Dechales}, \textit{Cursus}, tract.\:VII, lib.\;I, digr.\;IX %
(S.\;450b: einziges Diagramm).\cite{00124}%
}}%
}
\end{edfigure}%
\newpage%
\pstart%
\noindent%
\edtext{%
\textit{Afferuntur ex Svevia%
\indexO{Schwaben}
vitreae lacrymae%
\indexS{lacryma vitrea}
e vitro crassiori et fere viridi%
\indexS{vitrum viride}
confectae et solidae;
crassiores in medio},
sed \textit{desinunt in tenuissimum capillum.}%
\indexS{capillus tenuis}
\textit{Dominus Fieschi%
\indexN{\textls{Fieschi}, Sinibaldo, 1665 Inhaber einer Glasmanufaktur in Lyon}
qui anno 1665 Lugduni%
\indexO{Lyon}
fornacem%
\indexS{fornax}
ad crystallum%
\indexS{crystallus}
Veneta limpidiorem excitarat},
asseruit:
Materiam vitri hujus%
\indexS{materia vitri}
hoc habere peculiare,
quod
\textit{in aqua extincta
non abiret in friabilem pulverem%
\indexS{pulvis friabilis}
ut solet vitrum nostrum}%
\bracks{,}
\textit{quod} in pulverem \textit{abit,
si} ei \textit{immergatur dum incanduit;
dum haec materia colliquata in frigidam} incidit,
\textit{exteriores partes subito condensantur}
licet \textit{interioribus male compactis.}
Etiam \textit{vitrum commune praecipue crassius%
\indexS{vitrum commune}
nisi sensim frigefiat
externi aeris frigore diffringetur.}%
\indexS{aer externus}%
\indexS{frigus aeris}
Hinc \textit{artifices%
\indexS{artifex vasorum}
vas elaboratum%
\indexS{vas elaboratum ab artificibus}
rursus fornaci%
\indexS{fornax}
imponunt}
in
\edtext{%
quo \textit{recoquatur}%
\indexS{vas recoctum in fornace}%
}{%
\lemma{quo}%
\Bfootnote{%
\textit{(1)}~\angles{per}coquatur %
\textit{(2)}~\textit{recoquatur}%
~\textit{L}%
}}
et \textit{de} nimia \textit{rigiditate remittat.}%
\indexS{rigiditas vitri}
Eodem modo \textit{ferrum candens
frigidae
immersum}%
\bracks{,}%
\indexS{ferrum candens immersum}
\textit{tantum induit rigorem%
\indexS{rigor ferri}
ut facile frangi possit,
nisi denuo igni imponatur.}%
\indexS{ferrum impositum igni}%
}{%
\lemma{\textit{Afferuntur} \bracks{...} \textit{imponatur}}%
\Cfootnote{%
Vgl. a.a.O. %
(S.\;450b: Zitat mit Auslassungen).\cite{00124}
Ein nicht weiter ermittelter Sinibaldo Fieschi\indexN{\textls{Fieschi}, Sinibaldo, 1665 Inhaber einer Glasmanufaktur in Lyon} führte 1655 eine Glasmanufaktur in Lyon;\indexO{Lyon}
vgl. etwa \textsc{S.\;Fieschi}, Brief an J.-B.\;Colbert\indexN{\textls{Colbert}, Jean-Baptiste 1619\textendash1683} vom Februar 1665 (Bibliothèque Nationale de France, Département des manuscrits, Mélanges de Colbert 127, Bl.\;576\,r\textsuperscript{o}).\cite{02552}%
}}
\pend%
\begin{edfigure}[2.5em][0.0em]%
\includegraphics[scale=0.85]{%
gesamttex/edit_VIII-4/images/LH_35_09_17_005-006_d10_005v.pdf}
\edcaption[0.0em]{%
\lbrack\textit{Fig.\:10}\rbrack%
\edtext{}{%
\lemma{\lbrack\textit{Fig.\:10}\rbrack}%
\Cfootnote[nonum]{%
Vgl. \textsc{Dechales}, \textit{Cursus},
tract.\:VII, lib.\;I, digr.\;IX %
(S.\;453b u.\;454a: jeweils einziges Diagramm).\cite{00124}
Siehe hierzu und zur folgenden Beschreibung der Vakuumpumpe
(unten, \refpassage{LH_35_09_17_005-006_005v_jmt-1}{LH_35_09_17_005-006_005v_jmt-2}):
\textsc{Boyle}, \textit{New experiments}, Proemium u.\;Fig.\:1
(S.\,9--14 u.\;Tafel d.\;Abbildungen;\cite{00015} %
\textit{BW}~I, S.\;156, 160--162).\cite{00017}%
\indexN{\textls{Boyle} (Boylius, Boyl), Robert 1627\textendash1691}%
}}%
}
\end{edfigure}%
\newpage%
\pstart%
\noindent%
\edlabel{LH_35_09_17_005-006_005v_jmt-1}%
\edtext{%
\textit{AB} recipiens%
\indexS{recipiens capax 50 librarum aquae}%
\indexS{libra}
\edtext{%
\textit{50 librarum aquae capax.}%
}{%
\lemma{\textit{50 librarum aquae capax}}%
\Bfootnote{%
\textit{erg.\;L}%
}}
\textit{BC} epistomium cum verticillo \textit{D}.%
\indexS{epistomium cum verticillo}%
\indexS{verticillum}
\textit{In summitate} Recipientis
\textit{foramen 4 digitorum%
\indexS{foramen in summitate recipientis}
cum marginali labio
in conum turbinato%
\indexS{labium marginale turbinatum in conum}
ad altitudinem unius digiti,%
\indexS{digitus}
huic inseritur operculum
quod caemento ferruminatur%
\indexS{operculum recipientis ferruminatum caemento}
in cujus medio est foramen,%
\indexS{foramen instructum clavicula}
clavicula instructum,%
\indexS{clavicula operculi in summitate recipientis}
ut cum opus aer intromittatur.%
\indexS{aer intromissus recipienti}
Debet collum recipientis%
\indexS{collum recipientis ferruminatum caemento}
inseri tubo stanneo%
\indexS{tubus stanneus}
aut cupreo,%
\indexS{tubus cupreus}
et diligenter ferruminari caemento%
\indexS{caementum e pice resina et cineribus ligneis}
ex pice resina%
\indexS{pix resina}
et ligneis cineribus.%
\indexS{cinis ligneus}%
}%
}{%
\lemma{\textit{AB} recipiens \bracks{...} \textit{cineribus}}%
\Cfootnote{%
Vgl. \textsc{Dechales}, \textit{Cursus},
tract.\;VII, lib.\;I, digr.\;IX %
(S.\;453b: Zitat mit Auslassungen).\cite{00124}%
}}
\pend%
\pstart%
\edtext{%
\textit{%
Secunda pars machinae%
\indexS{machina (Boyleana)}
constat antlia \textit{E}%
\indexS{antlia constans cylindro cavo}
et ejus sustentaculo.%
\indexS{sustentaculum antliae}
Antlia cylindrum habet
cavitatis trium digitorum,%
\indexS{digitus}
huic aptatur Embolus%
\indexS{embolus antliae}
constans alio cylindro metallico minore et perforato,%
\indexS{cylinder metallicus perforatus}
cui imponitur segmentum corii
satis crassum et amplum,
optime lixivio praeparatum,%
\indexS{corium praeparatum lixivio}%
\indexS{lixivium}
ut cavitati cylindri coaptari possit,%
\indexS{cavitas cylindri}
et omnem aeri externo aditum praecludat.%
\indexS{aer externus}%
\indexS{aditus praeclusus aeri externo}
Manubrium autem emboli%
\indexS{manubrium emboli dentatum}
ferreum est %
et dentatum,
ita ut inseri possint denticuli rotae\;\textit{G}.}%
\indexS{denticulus rotae}%
\indexS{rota dentata}
Sed et
\textit{immittendum oleum intra cavitatem antliae,%
\indexS{oleum immittendum cavitati antliae}%
\indexS{cavitas antliae}%
\indexS{antlia constans cylindro cavo}
ut et facilius currat embolus,%
\indexS{embolus antliae}
et} melius \textit{excludatur aer.}%
\indexS{aer exclusus ex antlia}
Jam
\textit{obserata clavi \textit{CD},%
\indexS{clavis antliae}
embolum%
\indexS{embolus immissus summitati antliae}
ad summitatem usque}
\edtext{%
\textit{antliae} prius immissum%
\bracks{,}%
}{%
\lemma{\textit{antliae}}%
\Bfootnote{%
\textit{(1)}~\textit{immittas} %
\textit{(2)}~prius immissum%
~\textit{L}%
}}
\textit{quanta fieri poterit celeritate%
\indexS{embolus eductus ad imum antliae}
educes ad imum} usque \textit{antliae,%
\indexS{antlia constans cylindro cavo}
et versa clavicula}%
\indexS{clavicula antliae}
aditum aperies aeri%
\indexS{aditus aeris e recipiente in antliam}
\textit{ex recipiente in cavitatem} cylindri.
Jam clausa iterum clavi%
\indexS{clavis antliae}
\edtext{%
\textit{CD},
embolum%
\indexS{embolus antliae}%
}{%
\lemma{\textit{CD},}%
\Bfootnote{%
\bars{aerem \textit{streicht Hrsg.}} embolum %
~\textit{L}%
}}
rursus in cavitatem immittes%
\indexS{embolus immissus in cavitatem antliae}
aeremque expelles.%
\indexS{aer expulsus e cavitate antliae}%
\indexS{expulsio aeris e cavitate antliae}
Jamque rursus
educes.%
\indexS{embolus eductus e cavitate antliae}
Ita %
continuando%
\bracks{,}
semper tamen minus ac minus aeris extrahes.%
\indexS{aer extractus antlia}%
\indexS{extractio aeris antlia}%
}{%
\lemma{\textit{Secunda} \bracks{...} extrahes}%
\Cfootnote{%
Vgl. a.a.O. %
(S.\;454a: Zitat mit Auslassungen).\cite{00124}%
}}%
\edlabel{LH_35_09_17_005-006_005v_jmt-2}%
\pend%
\pstart%
\edtext{%
\textit{%
Arida agni vesica%
\indexS{vesica agni arida}%
\indexS{agnus}
capax duarum tertiarum pintae,%
\indexS{pinta}
per aeris inclusi recessum%
\indexS{recessus aeris inclusi}%
\indexS{aer inclusus vesicae}
amisit de suo pondere granum,%
\indexS{granum ponderis}
et octavam grani partem.
Vitreus recipiens aere exhaustus%
\indexS{recipiens exhaustus aere}
levior fuit uncia et tertia unciae parte,%
\indexS{uncia}
fuit autem capax mensurarum 32.%
\indexS{mensura}%
}%
}{%
\lemma{\textit{Arida} \bracks{...} \textit{mensurarum 32}}%
\Cfootnote{%
Vgl. a.a.O., experim.\;11 (S.\;455a);\cite{00124}
nach \textsc{Boyle}, \textit{New experiments}, Nr.\;1
(S.\;29\,f.;\cite{00015} \textit{BW}~I, S.\;167\,f.).\cite{00017}%
}}
\edtext{%
\textit{Aeolipila calefacta}%
\indexS{aeolipila calefacta}
occluso \textit{foramine}
ponderata est,
et post iterum
\textit{admisso
\edtext{%
aere,%
\indexS{aer admissus aeolipilae}
XI granis%
\indexS{granum ponderis}%
}{%
\lemma{aere,}%
\Bfootnote{%
\bars{et \textit{streicht Hrsg.}} \textit{XI granis}%
~\textit{L}%
}}
ponderosior fuit.}%
}{%
\lemma{\textit{Aeolipila} \bracks{...} \textit{fuit}}%
\Cfootnote{%
Vgl. \textsc{Dechales}, \textit{Cursus},
tract.\;VII, lib.\;I, digr.\;IX, experim.\;26
(S.\;457a: Zitat mit Auslassungen);\cite{00124}
nach \textsc{Boyle}, \textit{New experiments}, Nr.\;36
(S.\;286\,f.;\cite{00015} \textit{BW}~I, S.\;258\,f.).\cite{00017}%
}}
\edtext{%
\textit{Vesica aere plena%
\indexS{vesica plena aere}
inventa ponderosior XIV granis%
\indexS{granum ponderis}
se ipsa vacua,%
\indexS{vesica vacua}
erat autem capax 14 librarum aquae},%
\indexS{libra}
ergo
\indexS{proportio aquae ad aerem}%
\indexS{proportio aeris ad aquam}%
aqua ad aerem,
\textit{ut libra ad granum},%
\indexS{granum ponderis}
seu ut 7600 ad 1.
\ledpb%
\edtext{%
Ricciolus%
\indexN{\textls{Riccioli} (Ricciolus), Giovanni Battista 1598\textendash1671}%
}{%
\lemma{Ricciolus}%
\Cfootnote{%
Nicht ermittelt.
Siehe aber \textsc{Boyle}, \textit{New experiments}, N.\;36
(S.\;289;\cite{00015} \textit{BW}~I, S.\;259).\cite{00017}%
}}
vult ut 10000
ad 1.
Verum
\textit{ex Aeolipila%
\indexS{aeolipila calefacta}
alia proportio%
\indexS{proportio aquae ad aerem}%
\indexS{proportio aeris ad aquam}
nempe ut 1000 ad 1 circiter.
\edtext{%
Mersennus\indexN{\textls{Mersenne} (Mersennus), Marin 1588\textendash1648}%
}{%
\lemma{\textit{Mersennus}}%
\Cfootnote{%
Siehe \textsc{M.\;Mersenne}, \textit{Cogitata physico-mathematica}, De hydraulicis et pneumaticis phaenomenis, prop.\;29 (Paris 1644, S.\;141\,f.).\cite{00130}
Siehe zudem \textsc{Boyle}, \textit{New experiments}, N.\;36
(S.\;291;\cite{00015} \textit{BW}~I, S.\;260).\cite{00017}%
}}
ut 1300 ad 1.}%
\indexS{proportio aquae ad aerem}%
\indexS{proportio aeris ad aquam}%
}{%
\lemma{\textit{Vesica} \bracks{...} \textit{ad~1}}%
\Cfootnote{%
Vgl. \textsc{Dechales}, \textit{Cursus}, tract.\;VII, lib.\;I, %
digr.\;IX, experim.\;27 (S.\;457a: Zitat mit Auslassungen);\cite{00124}
nach \textsc{Boyle}, \textit{New experiments}, Nr.\;36
(S.\;288--291;\cite{00015} \textit{BW}~I, S.\;259\,f.).\cite{00017}%
}}
\pend%
\vspace{1.0em}%
\pstart%
\noindent%
\edlabel{LH_35_09_17_005-006_005v_dfg-1}%
\edtext{%
\textit{Medicorum cucurbitulae ita applicantur:%
\indexS{cucurbitula medicorum}
calefit aer inclusus quantum fieri potest,%
\indexS{aer inclusus cucurbitulae}%
\indexS{aer calefactus}
deinde carni applicantur firmiter,%
\indexS{cucurbitula applicata carni}
ita ut nullus aeri externo pateat aditus,%
\indexS{aditus patens aeri externo}
tunc frigefit cucurbitula%
\indexS{cucurbitula frigefacta}
vel infusa frigida,%
\indexS{frigida (aqua) infusa}
vel applicatione pannorum madidorum,%
\indexS{pannus madidus}
tunc caro elevatur intra cucurbitulam,%
\indexS{caro elevata intra cucurbitulam}
saepe etiam ossa,%
\indexS{os luxatum}
si quae sint luxata;
si incidatur caro,%
\indexS{caro incisa}
sanguis multus exilit.%
}%
\indexS{sanguis exiliens}%
\edlabel{LH_35_09_17_005-006_005v_dfg-2}%
}{%
\lemma{\textit{Medicorum} \bracks{...} \textit{exilit}}%
\Cfootnote{%
Vgl. \textsc{Dechales}, \textit{Cursus}, tract.\;VII, lib.\;I, %
digr.\;VI (S.\;442a: Zitat mit Auslassungen).\cite{00124}
Die Passage ist zweimal exzerpiert;
vgl. oben, \refpassage{LH_35_09_17_005-006_005r_ewk-1}{LH_35_09_17_005-006_005r_ewk-2}.%
}}
\pend%
\begin{edfigure}[0.00em][1.0em]
\includegraphics[scale=0.85]{%
gesamttex/edit_VIII-4/images/LH_35_09_17_005-006_d11_005v.pdf}
\edcaption[0.25em]{\lbrack\textit{Fig.\:11}\rbrack%
\edtext{}{%
\lemma{\lbrack\textit{Fig.\:11}\rbrack}%
\Cfootnote[nonum]{%
Vgl. \textsc{Dechales}, \textit{Cursus}, tract.\:VII, lib.\;I, %
digr.\:VII, n.\;16
(S.\;444a: einziges Diagramm).\cite{00124}%
}}%
}%
\end{edfigure}
\pstart%
\edlabel{LH_35_09_17_005-006_005v_ukp-1}%
\edtext{%
\textit{Sit tubus capitatus \textit{bc}%
\indexS{tubus capitatus}
in cujus summitate circumligetur frustulum suillae vesicae,%
\indexS{vesica suilla}
quod frustulum annectatur ipsi \textit{d} lanci bilancis \textit{de},%
\indexS{lanx bilancis}%
\indexS{bilanx}
impletoque tubo \textit{bc}
fiat experimentum more consueto}%
\indexS{mos consuetus}%
\lbrack:\rbrack\
\textit{vesica deberet intra caput tubi praecipitari,%
\indexS{caput tubi}%
\indexS{vesica praecipitata intra tubum}
quod facile impedietur
adjectis in adversa lance ponderibus}%
\indexS{pondus adjectum in lance adversa}%
\lbrack;\rbrack\
\textit{fiat igitur aequilibrium%
\indexS{aequilibrium in bilance}
ita ut detracto vel minimo pondere%
\indexS{pondus vel minimum detractum}
praecipitanda esset vesica}%
\indexS{vesica praecipitata intra tubum}%
\lbrack;\rbrack\
\textit{si supponas digitum tubo,
supponasque Mercurium%
\indexS{mercurius contentus in tubo}
in eo contentum,
invenies illud}
(+\:an illum\:+)
\textit{aequiponderare ponderibus%
\indexS{mercurius aequiponderans in bilance}
in adversa lance positis.}%
\edlabel{LH_35_09_17_005-006_005v_ukp-2}%
}{%
\lemma{\textit{Sit tubus} \lbrack...\rbrack\ \textit{positis}}%
\Cfootnote{%
Vgl. a.a.O.\cite{00124}
(S.\;444a: Zitat mit Auslassungen).
Die Passage setzt einen früheren Auszug fort;
vgl. oben, \refpassage{LH_35_09_17_005-006_005r_svj-1}{LH_35_09_17_005-006_005r_svj-2}.%
}}
\pend%
\pstart%
\edtext{%
Aerem includi Mercurio%
\indexS{aer inclusus mercurio}
multis indiciis constat.%
\indexS{indicium}
Sane aerem esse in liquidis,%
\indexS{aer in liquidis}%
\indexS{liquidum}
patet ex \textit{dissolutionibus}%
\indexS{dissolutio per aquas stygias}
quae fiunt \textit{per aquas stygias},%
\indexS{aqua stygia}
item \textit{in machina}%
\indexS{machina ministrans ventum}
qua \textit{aqua praecipitatur}%
\indexS{aqua praecipitata}
ut \textit{follibus%
\indexS{follis}
ventum ministret}.%
\indexS{machina ministrans ventum}%
\indexS{ventus ministratus per machinam}
Mercurius \textit{est maxime volatilis};%
\indexS{mercurius volatilis}%
\indexS{volatilitas mercurii}
\edlabel{LH_35_09_17_005-006_005v_jkr-1}%
ajunt
\textit{qui} eum \textit{in officinis servant},%
\indexS{mercurius servatus in officinis}%
\indexS{officina}
\textit{in singulos}
\edtext{%
\textit{anno}\bracks{\textit{s}}%
\indexS{annus singulus}%
}{%
\lemma{\textit{anno}}%
\Bfootnote{%
\textit{L\;ändert Hrsg. nach Vorlage}%
}}
\textit{20\textsuperscript{ma} aut}
\edtext{%
\bracks{\textit{25\textsuperscript{ta}}}%
}{%
\lemma{\textit{25\textsuperscript{ma}}}%
\Bfootnote{%
\textit{L\;ändert Hrsg.}%
}}
\textit{parte decrescere},%
\indexS{decrementum mercurii servati}
\ledpb%
quacunque diligentia%
\indexS{diligentia in mercurio servando}
servetur.%
\edlabel{LH_35_09_17_005-006_005v_jkr-2}
\edtext{%
Si eodem
\edtext{%
Mercurio%
}{%
\lemma{Mercurio}%
\Bfootnote{%
\textit{erg.\;L}%
}}
experimentum Torricellianum%
\indexS{experimentum Torricellianum}
repetas%
\bracks{,}
\textls{\textit{non tam descendit}}%
\bracks{;}%
}{%
\lemma{\textit{Am Rand:}}%
\Afootnote{%
\Denarius%
}}
\textit{post iterata saepe experimenta%
\indexS{experimentum iteratum}
multo fit deterior,%
\indexS{mercurius deterioratus}
et mobilitate,%
\indexS{mobilitas mercurii}
splendore}%
\indexS{splendor mercurii}
etc. imminuitur.%
}{%
\lemma{Aerem \bracks{...} imminuitur}%
\Cfootnote{%
Vgl. \textsc{Dechales}, \textit{Cursus}, tract.\:VII, lib.\;I, %
digr.\:VII (S.\;445a--b: Zitat mit Auslassungen).\cite{00124}
Die Passage ist zum Teil zweimal exzerpiert;
vgl. oben, \refpassage{LH_35_09_17_005-006_005r_scs-1}{LH_35_09_17_005-006_005r_scs-2}.%
}}
\bracks{%
\href{\Scan{LH_35_09_17_006r}}%
{6\,r\textsuperscript{o}}%
}
\pend%
\vspace{1.0em}%
\pstart%
\noindent%
\edtext{%
\textit{Si capillaris tubus%
\indexS{tubus capillaris immersus}
aquae immergatur,
in eo ascendet aqua ad 5 digitos} usque.%
\indexS{aqua ascendens in tubo capillari}
Hinc \textit{uni\-ver\-sa\-le principium%
\indexS{principium universale}
quo liquor se insinuat et ascendit%
\indexS{liquor ascendens per arenas}\indexS{liquor ascendens per cineres}\indexS{liquor ascendens per fibras arborum}
per arborum fibras,%
\indexS{fibra arboris}%
\indexS{arbor}
per medios cineres,%
\indexS{cinis}
arenas,%
\indexS{arena}
aliaque.}%
}{%
\lemma{\textit{Si capillaris} \bracks{...} \textit{aliaque}}%
\Cfootnote{%
Vgl. \textsc{Dechales}, \textit{Cursus}, tract.\:VII, lib.\;I, %
digr.\;IX, experim.\;25 (S.\;456b: Zitat mit Auslassungen).\cite{00124}%
}}
\pend%
\pstart%
\edtext{%
Difficultas vulgo,%
\indexS{difficultas vulgo}
si aqua in aquam
\edtext{subjectam}{%
\lemma{subjectam}%
\Bfootnote{%
\textit{erg.\;L}%
}}
gravitat%
\indexS{aqua gravitans in aquam subjectam}%
\indexS{aqua subjecta}
quomodo fit,
ut in fundo%
\indexS{fundus}
pondus aquae vix sentiamus?%
\indexS{pondus aquae in fundo}%
}{%
\lemma{Difficultas \bracks{...} sentiamus}%
\Cfootnote{%
Vgl. a.a.O., %
digr.\;X, obiect.\;4 (S.\;460a).\cite{00124}%
}}
\edtext{%
Equidem \textit{in acervo}%
\indexS{acervum ovorum}%
\indexS{ovum}
\edtext{%
\textit{ovorum inferiora}%
}{%
\lemma{\textit{ovorum}}%
\Bfootnote{%
\textit{(1)}~acervis ovor %
\textit{(2)}~\textit{inferiora}%
~\textit{L}%
}}
\textit{non franguntur a} super\textit{positorum pondere},
itaque \textit{si homo manum ovis supponeret,%
\indexS{manus supposita ovis}
non majorem sentiret}
\edtext{%
\textit{pressionem%
\indexS{pressio ovorum in manum suppositam}
quam} unum \textit{ovum},
quae%
}{%
\lemma{\textit{pressionem quam}}%
\Bfootnote{%
\textit{(1)}~unum \textit{ovum} qu %
\textit{(2)}~unum \textit{ovum}, quae%
~\textit{L}%
}}
minor est
quam ut ovum frangat.%
\indexS{pressio frangens ovum}%
\indexS{ovum fractum pressione}
Si \textit{manus acervo milii supponatur}%
\indexS{acervum milii}%
\indexS{manus supposita acervo milii}
\textit{non sentitur pondus milii}%
\indexS{pondus milii}%
\indexS{milium}
(\textit{in arenae cumulo}%
\indexS{cumulus arenae}%
\indexS{arena}
\edtext{%
paulo%
}{%
\lemma{paulo}%
\Bfootnote{%
\textit{erg.\;L}%
}}
majus \textit{sentitur pondus}).%
\indexS{pondus cumuli arenae}
In acervo milii%
\indexS{acervum milii}%
\indexS{milium}
\textit{distribuitur pondus singulis partibus},%
\indexS{pondus distributum singulis partibus}
hinc \textit{nulla partium separatio}
seu \textit{laesio};%
\indexS{laesio}
\textit{partes se invicem sustentant}%
\indexS{partes sustentantes se instar fornicis}
instar \textit{fornicis};%
\indexS{fornix}
\textit{divisibilitas facilis},%
\indexS{partes divisibiles}%
\indexS{divisibilitas partium}
quia \textit{milii grana laevia et polita}%
\indexS{granum milii}%
\indexS{milium}
\edtext{%
faci\bracks{un}t,%
}{%
\lemma{facit}%
\Bfootnote{%
\textit{L\;ändert Hrsg.}%
}}
ut facile sub acervo moveatur manus.%
\indexS{manus mota sub acervo milii}
\edtext{%
Si quae%
}{%
\lemma{Si}%
\Bfootnote{%
\textit{(1)}~aliae %
\textit{(2)}~quae%
~\textit{L}%
}}
\textit{partes milii manum premunt},
\textit{aliae} %
\textit{ad latera sustentant,
unde vix dignosci potest
in quam partem impellatur},
et \textit{sentitur potius mollis compressio%
\indexS{compressio mollis}
quam gravitatio.%
\indexS{gravitatio acervi milii}%
}%
}{%
\lemma{Equidem \bracks{...} \textit{gravitatio}}%
\Cfootnote{%
Vgl. a.a.O. %
(S.\;460b: Zitat mit Auslassungen).\cite{00124}%
}}
\edtext{%
At \textit{aperto foramine sub manu}%
\indexS{foramen apertum sub manu}
\edtext{%
manus magis%
}{%
\lemma{manus}%
\Bfootnote{%
\textit{(1)}~melius %
\textit{(2)}~magis%
~\textit{L}%
}}
\textit{percipiet pondus}.%
\indexS{manus percipiens pondus}%
\indexS{pondus perceptum manu}%
}{%
\lemma{At \textit{aperto} \bracks{...} \textit{pondus}}%
\Cfootnote{%
Vgl. a.a.O. %
(S.\;461a: Zitat mit Auslassungen).\cite{00124}%
}}
\pend%
\pstart%
\edtext{%
Quidam ridicule putabat
\textit{in malacia}%
\indexS{malacia}
homines%
\indexS{homo intra navem}
qui sunt intra navem%
\indexS{navis in malacia}
debere \textit{alligare funem ad puppim},%
\indexS{funis alligatus puppi}%
\indexS{puppis}
et ita \textit{trahere ad proram}%
\indexS{prora}
atque ita movere navem.%
\indexS{navis in malacia}%
}{%
\lemma{Quidam \bracks{...} navem}%
\Cfootnote{%
Vgl. a.a.O., %
digr.\;XI (S.\;462b: Zitat mit Auslassungen).\cite{00124}%
}}
\edtext{%
Quidam studiosus%
\indexS{studiosus}
\textit{lanternam componere} volebat,%
\indexS{lanterna mota impetu fumi}
\textit{quae solius fumi%
\indexS{fumus erumpens e candela}
\edtext{%
ex candela%
\indexS{candela infixa lanternae}%
}{%
\lemma{\textit{ex}}%
\Bfootnote{%
\textit{(1)}~lanterna %
\textit{(2)}~\textit{candela}%
~\textit{L}%
}}
erumpentis impetu}%
\indexS{impetus fumi erumpentis}
moveretur%
\bracks{:}
cum igitur \textit{alas}%
\indexS{ala lanternae}
prout oportet
\textit{detorsisset}
ut in machinis aliis%
\indexS{machina circumacta fumo}
quae fumo circumaguntur,%
\indexS{fumus circumagens machinam}
\textit{lanterna}%
\indexS{lanterna mota impetu fumi}
tamen
\textit{licet duabus candelis%
\indexS{candela infixa lanternae}
adhibitis}
non \textit{nutabat}.
\ledpb%
Sed causa erat,%
\indexS{causa immobilitatis}
quod candela%
\indexS{candela infixa lanternae}
esset lanternae%
\indexS{lanterna mota impetu fumi}
infixa
et deberet simul moveri;
cum vero candela esset immota,
et lanterna circa eam moveretur,
res successit
et \textit{machina} magno \textit{impetu mota} est.%
\indexS{machina mota impetu magno}
Eodem
\edtext{%
mod\bracks{o}%
}{%
\lemma{modus}%
\Bfootnote{%
\textit{L\;ändert Hrsg.}%
}}
ventus ex terra quippe assurgens,%
\indexS{ventus assurgens e terra}%
\indexS{terra}
ictus lapidis%
\indexS{ictus lapidis recidentis in terram}%
\indexS{lapis recidens in terram}
in terram recidentis,%
\indexS{terra}
aliaque id genus
non magis terram movent
quam homines in nave navim,%
\indexS{navis in malacia}%
\indexS{homo intra navem}
si in nihil extra navim impingant.
Nam et aer a terra sustinetur.%
\indexS{aer sustentus a terra}%
\indexS{terra sustinens aerem}%
}{%
\lemma{Quidam \bracks{...} sustinetur}%
\Cfootnote{%
Vgl. a.a.O. %
(S.\;463a: Zitat mit Auslassungen).\cite{00124}%
}}
\pend%
\vspace{2.5em}%
\pstart%
\noindent%
\centering%
*     *     *
\pend%
\vspace{3.0em}%
\selectlanguage{german}%
\frenchspacing%
\footnotesize%
\pstart%
\noindent%
\textls{Ausführliche Datierungsgründe:}%
\pagelabel{41182_dat}
Das vorliegende undatierte Stück N.\;\ref{41182-1} umfasst Leibnizens ausführliche Auszüge mit Bemerkungen aus \textsc{C.\:F.\:M.\;Dechales}, \textit{Cursus seu mundus mathematicus} (3~Bde, Lyon 1674), tract.\;VII: \textit{Statica}, lib.\;I: \textit{Digressiones physicae} (Bd.\;I, S.\;433--464).\cite{00124}
Die Auszüge sind auf einem Papierbogen überliefert (LH~XXXV\;9,\;17 Bl.\;5--6), der weder Spuren eines Wasserzeichens noch sonstige datierungsrelevante äußere Merkmale aufweist.
Derselbe Bogen überliefert auch die wohl zum gleichen Zeitpunkt entstandenen Auszüge aus Dechales'\indexN{\textls{Dechales} (Chalesius), Claude François Milliet 1621\textendash1678} \textit{Cursus}, tract.\;VII, lib.\;III: \textit{De descensu gravium in planis inclinatis, et funependulis} (Bd.\;I, S.\;489--524, bes. S.\;491--496), die in diesem Band unter N.\;\ref{41182-2} ediert sind.
Beide Stücke N.\;\ref{41182-1} und N.\;\ref{41182-2} gehören zu einer umfangreicheren Gruppe von Exzerpten aus Dechales'\indexN{\textls{Dechales} (Chalesius), Claude François Milliet 1621\textendash1678} \textit{Cursus},\cite{00124} die nach heutigem Wissensstand folgende Manuskripte umfasst:
LH~XXXV\;9,\;17 Bl.\;1--4; %
LH~XXXVI Bl.\;54--57, 60--61; %
LH~XXXVII\;5 Bl.\;2--3 %
(sämtliche werden voraussichtlich in künftigen Bänden von \textit{LSB}\;VIII veröffentlicht).
Darunter weisen drei Bogen (Bl.\;1--4; 60--61) kein Wasserzeichen auf, wohingegen drei Bogen (Bl.\;2--3; 54--57) mit ein und derselben Marke versehen sind.
Dieses Wasserzeichen kennzeichnet eine Papiersorte, die Leibniz im Zeitraum 1683--1685/6, vermutlich während seiner Aufenthalte im Harz,\indexO{Harz} verwendete.
Innerhalb dieses Zeitraumes dürfte die ganze Gruppe der Handschriften entstanden sein, die Auszüge aus Dechales' \textit{Cursus} überliefern: darunter wohl auch das vorliegende Stück N.\;\ref{41182-1}.
In dieser Hinsicht ist besonders relevant, dass das besagte Wasserzeichen auch im Bogen LH~XXXV\;9,\;17 Bl.\;2--3 nachweisbar ist, welcher Auszüge aus dem Traktat VI (\textit{Mechanica}) überliefert, der in Dechales' \textit{Cursus}\cite{00124} dem Traktat VII (\textit{Statica}) unmittelbar vorhergeht.
Die Annahme liegt nahe, dass Leibniz beide Traktate nacheinander exzerpiert hat.
\pend%
\pstart%
Weiter eingrenzen lässt sich die Datierung der vorliegenden Auszüge N.\;\ref{41182-1} unter Berufung auf folgenden Umstand:
In Leibnizens physikalischer Abhandlung \textit{Cogitationes novae quomodo formetur sonus et per aerem propagetur, atque in organo auditus exprimatur} findet sich ein ausführliches, zum Teil wörtliches Zitat aus \textsc{Dechales}, \textit{Cursus seu mundus mathematicus}, tract.\;XXII, prop.\;16 (Bd.\;III, S.\;24a--b);\cite{00124} vgl. \textit{LSB} VIII,\;3 N.\;12\textsubscript{3}, S.\;135.9--136.3.\cite{02529}
Zu dem Zeitpunkt, als diese Abhandlung entstand, hätte sich Leibniz offenbar gründlich mit Dechales' Werk, oder wenigstens mit Teilen davon, auseinandergesetzt.
In der ersten Fassung der \textit{Cogitationes novae} -- d.\,h. dem zwischen Frühjahr 1682 und Winter 1682/1683 verfassten Konzept \textit{LSB} VIII,\;3 N.\;12\textsubscript{3}, \textit{L\textsuperscript{1}} -- fehlte jedoch das besagte Zitat aus Dechales'\indexN{\textls{Dechales} (Chalesius), Claude François Milliet 1621\textendash1678} \textit{Cursus}\cite{00124} noch gänzlich.
\ledpb%
Dieses wurde (wie die gesamte Passage \textit{LSB} VIII,\;3 N.\;12\textsubscript{3}, S.\;135.9--137.4) erst zu dem Zeitpunkt hinzugefügt, als Leibniz die von seinem Sekretär Jobst Dietrich Brandshagen\indexN{\textls{Brandshagen} (Brondhaguen), Jobst Dietrich 1659\textendash nach 1716} angefertigte Abschrift \textit{l} von \textit{L\textsuperscript{1}} überarbeitete, d.\,h. zwischen der zweiten Hälfte Juni 1684 und dem Spätfrühling 1685 (siehe zur Datierung der verschiedenen Textschichten der \textit{Cogitationes novae}\cite{02529} die editorische Erläuterung in \textit{LSB} VIII,\;3 N.\;12, bes. S.\;91.19--92.40).\cite{02526}
Zu diesem Zeipunkt hatte sich Leibniz mit großer Wahrscheinlichkeit nicht nur mit dem dritten Band von Dechales'\indexN{\textls{Dechales} (Chalesius), Claude François Milliet 1621\textendash1678} Werk befasst, sondern auch mit den oben erwähnten Traktaten im ersten Band des \textit{Cursus}.\cite{00124}
Daraus dürften auch die vorliegenden Auszüge N.\;\ref{41182-1} entstanden sein.
Diese sollten demgemäß nicht vor dem Winter 1682/1683 und nicht nach der ersten Hälfte 1685 verfasst worden sein.
Eine etwas frühere oder eine etwas spätere Datierung sind zwar nicht ausgeschlossen, aber unwahrscheinlich.
\pend%
\selectlanguage{latin}%
\frenchspacing%
\normalsize%
